%! Equivalent Representations of Polynomial and Rational Expressions — Unit 3, Lesson 3.5 — Homework
\documentclass[10pt]{article}
\usepackage{precalculus-article}
\usepackage{precalculus-boxes}
\newcommand{\TallMath}[1]{$\displaystyle #1\rule[-1.4em]{0pt}{3.2em}$}

\begin{document}

\pageheader{Unit 3, Lesson 3.5}{Homework}
\namedateperiod

\vspace{0.12in}

\begin{enumerate}[itemsep=18pt]

\item For each rational function, state which form (\textbf{factored} or \textbf{standard}) is
      more useful for the requested task, then answer the question.

      \begin{enumerate}[label=(\alph*), itemsep=10pt]
        \item $r(x) = \dfrac{x^2+3x-10}{x^2-4}$. \quad Find all zeros, holes, and VAs.

              \vspace{0.06in}
              Best form: \blank{3cm}

              Factor numerator: \blank{5cm} \quad Factor denominator: \blank{5cm}

              Zeros: \blank{3cm} \quad Holes: \blank{3cm} \quad VAs: \blank{3cm}

        \item $s(x) = \dfrac{3x^5 - 2x^2 + 1}{4x^3 - x}$. \quad Describe the end behavior.

              \vspace{0.06in}
              Best form: \blank{3cm}

              As $x \to +\infty$: $s(x) \to$ \blank{3cm} \quad As $x \to -\infty$: $s(x) \to$ \blank{3cm}
      \end{enumerate}

\item Perform polynomial long division. Show all steps.

      Divide $f(x) = 3x^3 - 5x^2 + 2x - 4$ by $g(x) = x - 2$.

      \vspace{0.1in}
      Work space:

      \vspace{1.0in}

      $q(x) = $ \blank{5cm} \quad $r = $ \blank{2cm}

      Write as $f(x) = g(x)\cdot q(x) + r$: \blank{8cm}

\item A rational function is defined by $\displaystyle R(x) = \dfrac{x^2 + x - 6}{x - 1}$.

      \begin{enumerate}[label=(\alph*), itemsep=8pt]
        \item Perform polynomial long division.

              \vspace{0.06in}
              $R(x) = $ \blank{8cm}

        \item State the equation of the slant asymptote. Explain in one sentence why this is
              the slant asymptote.

              \vspace{0.06in}
              Slant asymptote: $y = $ \blank{3cm}

              \writeline

        \item What is the remainder when $x^2+x-6$ is divided by $x-1$? What does this tell
              you about $R(1)$?

              \vspace{0.06in}
              Remainder: \blank{2cm} \quad $R(1) = $ \blank{2cm}
      \end{enumerate}

\item Use Pascal's Triangle to expand each expression. Show the row of Pascal's Triangle you used.

      \begin{enumerate}[label=(\alph*), itemsep=10pt]
        \item $(x + 3)^3$

              \vspace{0.06in}
              Row used: \blank{5cm}

              $(x+3)^3 = $ \blank{7cm}

        \item $(2x - 1)^4$ \quad \textit{(Hint: let $a = 2x$ and $b = -1$)}

              \vspace{0.06in}
              Row used: \blank{5cm}

              $(2x-1)^4 = $ \blank{9cm}
      \end{enumerate}

\item \textbf{(AP-Style Multiple Choice)} Which of the following statements about
      $\displaystyle p(x) = \dfrac{x^3 - 8}{x - 2}$ is correct?

      \medskip
      \begin{enumerate}[label=(\Alph*), itemsep=4pt]
        \item $p$ has a vertical asymptote at $x = 2$.
        \item $p$ has a hole at $x = 2$ with $y$-coordinate $12$.
        \item $p$ simplifies to $x^2 - 2x + 4$ for all real $x$.
        \item $p$ has a slant asymptote of $y = x^2 + 2x + 4$.
      \end{enumerate}

      \vspace{0.06in}
      Answer: \blank{1.5cm} \quad Show your reasoning:

      \vspace{0.06in}
      \writelines{3}

\end{enumerate}

\end{document}
